\section{模型检验与稳健性分析}
\label{sec:validation}

前述模型连续运行数百天，并在一天内多次更新决策。对于这类时序优化，求解器给出“最优解”只说明数学程序完成了优化，并不能自动保证方案满足真实运行逻辑。一个可信的调度结果至少应满足三点：每个时段的电量能够守恒，储能状态能够连续传递，且任一决策只能使用当时已经获得的信息。基于这一原则，下面不再重复模型求解过程，而是从物理轨迹、时间因果、结算逻辑和关键参数四个方面对最终结果进行独立复核。

\subsection{物理可行性与数值一致性}

首先把最终输出逐时重新代入能量平衡式和储能状态递推式。若这一层不能通过，即使全年费用看起来合理，也可能只是单位换算、储能递推或约束遗漏造成的“数值低成本”，并检查储电量是否始终落在允许区间、主方案是否出现同时充放电，以及由逐时交易重新汇总出的费用能否回到论文报告值。结果列于表~\ref{tab:validation}。

\begin{table}[htbp]
    \caption{物理残差与费用回代误差均低于设定校验阈值}
    \label{tab:validation}
    \centering
    \zihao{-5}
    \begin{tabular}{l l l l}
        \toprule
        检验项目 & 最大偏差或范围 & 判定标准 & 结论 \\
        \midrule
        能量平衡残差 & $<10^{-12}$ kWh & $10^{-7}$ kWh & 通过 \\
        储电量递推残差 & $<10^{-12}$ kWh & $10^{-7}$ kWh & 通过 \\
        储电量取值范围 & $[1200,\ 10800]$ kWh & 题设区间 & 通过 \\
        主方案充放电重叠量 & $<10^{-8}$ kWh & $10^{-5}$ kWh & 通过 \\
        问题三费用回代差 & $<10^{-8}$ 元 & $10^{-7}$ 元 & 通过 \\
        结果工作簿回读差异 & $<10^{-10}$ & 数值一致 & 通过 \\
        \bottomrule
    \end{tabular}
\end{table}

表~\ref{tab:validation} 中两类状态方程的残差均远低于数值容差，表中也给出求解结果的误差量级与费用回代精度。说明每一度电在模型中都有明确来源和去向，储能也没有通过越界或异常耗散获得额外收益。加入充放电互斥约束后，充、放电重叠量只剩浮点舍入量级。按逐时计划购电、调整交易和紧急购电重新结算，所得总费用与目标函数一致；五个结果工作簿重新读入后，也与求解明细逐项对应。后文讨论的费用差异可以解释为信息和策略差异，基本物理约束与文件写入均已通过独立校验。

\subsection{时间因果性的反事实扰动检验}

物理可行仍不足以证明计划真实可执行。最容易被忽略的风险是“未来信息穿越”：如果 $0{:}00$ 的全天计划会随着当天尚未发生的真实光伏或负载而变化，那么模型实际上提前知道了未来，即使费用更低也失去实际意义。为直接检查这一点，本文采用反事实扰动：固定某个决策节点之前的全部历史信息，只改变节点之后才会实现的数据，再观察当前计划是否随之改变。

在问题二中，分别修改当天未来真实负荷与真实光伏，$0{:}00$ 计划保持不变；在问题三中，改变 $6{:}00$ 之后的实际源荷，不会反向改变 $6{:}00$ 已经形成的计划；在问题四的因果价格扩展中，未来实际电价的扰动同样不能影响扰动发生前的决策。所有非预见性检验的最大计划差异均低于 $10^{-10}$。作为对照，在直接使用附件4完整价格的题面主情形中，改变未来价格会相应改变计划，说明前述“不变”确实来自信息隔离；对照情形也证明扰动本身能够影响计划。

滚动模型还需要保证计划能够正确延续。仅启用 $0{:}00$ 和 $6{:}00$ 两个更新节点时，$12{:}00$ 以后应继续沿用 $6{:}00$ 已经生效的未来方案，而不能因为 $12{:}00$ 节点被关闭就退回日初计划。逐段追踪表明，关闭节点后最近一次有效方案得到保留，且每个执行块都从上一阶段的实测储电量继续递推。由此，问题三各更新节点之间的费用差才真正对应新增信息，而不会混入错误的计划回退效应。

\subsection{调整结算与滚动计划的一致性检验}

问题三中，调整费用会改变模型愿意上调还是下调购电，因此费用公式必须和优化过程保持一致。先用单时段手算核对主模型与对照方案。设正常电价为 $p$，当计划购电量 $G^p=5000$ kWh、执行前调整为 $G^a=4000$ kWh 时，主模型按“未调整部分正常结算、少购部分支付违约费用”计算为 $4000p+0.5\times1000p=4500p$；对照方案把原计划费与违约费用并列相加，得到 $5000p+0.5\times1000p=5500p$。当计划由 $4000$ kWh 上调至 $5000$ kWh 时，两种计费方式均为 $4000p+1.5\times1000p=5500p$。两种计费方式只在实际购电量低于计划时产生分歧。对照方案下，少买 $1000$ kWh 会比不调整多付 $500p$，从而抑制向下修正；主模型下少买能够降低支出，但要为少买部分支付 $50\%$ 的违约费用，这与题面“违约电价”的措辞和经济含义一致。

不同结算规则会反过来改变最优调度。两种计费方式均从计划层重新求解，避免在同一条购电和储能轨迹上事后更换计价公式。随机源荷与价格序列上的交叉复算显示，每种计费方式下重新优化得到的方案，都不会被另一种方式的最优轨迹直接替代而获得更低费用。这说明费用规则已经真正进入决策。

对于多次滚动更新，同一未来时段可能先后出现在 $0{:}00$、$6{:}00$、$12{:}00$ 和 $18{:}00$ 的规划中，但实际执行时只采用交易前最后一次生效的方案。该处理与前述计划持续性检验共同保证：模型可以反复“看未来”，但不会重复执行多份计划，也不会在节点缺失时凭空改变既有安排。

\subsection{历史配置的复算一致性}

为确认正文调整未改变既有计算链，本文保留历史配置进行独立复算，并逐项比较旧计算与当前程序在相同参数下的输出。表~\ref{tab:historical-regression} 只报告一致性结论，不在正文重复列出旧版本费用数值，以避免历史结果与当前主方案混淆。

\begin{table}[htbp]
    \caption{相同参数下五组历史配置均可复算到一致结果}
    \label{tab:historical-regression}
    \centering
    \zihao{-5}
    \begin{tabular}{l c c}
        \toprule
        复算对象 & 旧计算与当前复算 & 最大绝对偏差/元 \\
        \midrule
        问题一 & 一致 & $<10^{-3}$ \\
        问题二 & 一致 & $<10^{-3}$ \\
        问题三 & 一致 & $<10^{-3}$ \\
        问题四-2 & 一致 & $<10^{-3}$ \\
        问题四-3 & 一致 & $<10^{-3}$ \\
        \bottomrule
    \end{tabular}
\end{table}

该复算说明，新增的信息条件、计费方式与储能约束并未破坏既有配置下的计算行为；历史结果可作为信息更充分或规则不同条件下的基准，而不与当前主方案混合比较。

\subsection{终端储电量策略的敏感性}
\label{sec:terminal}

问题一明确要求日初与日末储电量相同，而问题二、三没有给出同样的硬约束。如果完全取消终端要求，单日优化会倾向于在午夜前尽量放空电池，把本应留给下一日的储能价值转化为当前日的低费用。这种做法在单日目标函数上看似更优，却可能只是把成本推迟到下一日。为判断日循环假设是否过强，本文保留无终端约束方案作为对照。

两种终端策略的费用与储电量变化如图~\ref{fig:sensitivity} 所示。日循环策略全年费用为 \SensTerminalCycle\ 元；取消终端约束后费用降至 \SensTerminalFree\ 元，变化幅度仅为 \SensTerminalPct\%，但全年最终储电量下降到 \SensTerminalFreeSoc\ kWh，日末平均储电量也只有 \SensTerminalFreeMean\ kWh，已经接近 $1200$ kWh 的安全下限。

\begin{figure}[htbp]
    \centering
    \includegraphics[width=0.98\textwidth]{figures/sensitivity.pdf}
    \caption{取消终端约束仅小幅降费，却使日末储电量逼近安全下限}
    \label{fig:sensitivity}
\end{figure}

图~\ref{fig:sensitivity} 说明，无终端约束获得的少量费用优势主要来自持续压低日末储电量——日末状态大量集中在下限附近，并不意味着光伏或电价利用效率真正提高。其日末 SOC 标准差虽然为 \SensTerminalFreeStd\ kWh，小于日循环策略的 \SensTerminalCycleStd\ kWh，但这是因为大量状态集中在下限附近，不能解释为运行更稳健。考虑全年连续运行时下一日仍需要可比的储能起点，本文保留日循环作为主策略，并把无终端约束结果用于说明该假设对总费用的影响有限。

\subsection{储能效率取值的敏感性}
\label{sec:efficiency}

题目给出的“充放电效率为 $90\%$”可以有两种理解：其一是充电和放电效率分别为 $0.9$，对应往返效率 $0.81$；其二是整个充放电过程的往返效率为 $0.9$，若两侧对称，则可取 $\eta_c=\eta_d=\sqrt{0.9}$。由于这一解释会直接影响每次跨时段搬移电量的损耗，有必要在保持其他设置不变的条件下重新求解。结果见表~\ref{tab:efficiency}。

\begin{table}[htbp]
    \caption{提高往返效率会降低全年费用，但不改变储能与滚动更新的作用方向}
    \label{tab:efficiency}
    \centering
    \zihao{-5}
    \begin{tabular}{l r r r}
        \toprule
        效率定义 & 全年总费用/元 & 储能吞吐量/kWh & 紧急购电量/kWh \\
        \midrule
        定义A：$\eta_c=\eta_d=0.90$（往返 $0.81$） & \SensEffEOneCost & \SensEffEOneThroughput & \SensEffEOneEm \\
        定义B：$\eta_c=\eta_d=\sqrt{0.9}$（往返 $0.90$） & \SensEffETwoCost & \SensEffETwoThroughput & \SensEffETwoEm \\
        \bottomrule
    \end{tabular}
\end{table}

提高往返效率后，单位充入电量可以在后续释放更多有效电能，储能因而更愿意承担峰谷转移和偏差缓冲，全年费用随之下降。但两种效率定义下，储能的作用方向、风险购电的必要性以及滚动更新的相对优势均未改变。也就是说，效率解释会影响费用的绝对数值，却不会推翻本文关于“储能降低高价购电、更新信息降低偏差成本”的主要判断。主模型取 $\eta_c=\eta_d=0.9$，相当于采用较保守的效率解释。

\subsection{误差来源与适用边界}

模型剩余误差主要来自三类信息近似。其一，问题二的负载和光伏只根据历史序列预测，突发天气、节假日或异常用电会使近期历史失去代表性；其二，附件3只给出整点预报，将其插值到 $10$ min 能恢复调度时间尺度，却不能生成原预报中不存在的短时波动；其三，问题四主结果按题意直接使用附件4完整价格，而真实市场中未来价格的发布方式可能更受限制。本文分别通过滚动误差更新、储能实时反馈、节点对照和因果价格扩展量化了这些影响。

本文结果首先回答的是题目给定数据和交易规则下的调度问题。若向真实长期运行推广，还应引入天气预报、工作日特征、需求侧响应、储能衰减和价格发布机制等外生信息。上述限制不会破坏当前结果的内部一致性，但决定了模型从竞赛数据走向实际系统时还需要哪些新增信息与设备模型。
